% !TEX root = ../main.tex
\section{Implementación}

\subsection{Arquitectura del monorepo}

El proyecto se organiza como monorepo con \texttt{pnpm workspaces} y \texttt{Turborepo}. La tabla~\ref{tab:arquitectura} describe las cuatro capas y su responsabilidad funcional. La aplicación web utiliza Next.js App Router con TypeScript y Tailwind CSS, componentes estilo shadcn/ui, React Flow para el grafo, Recharts para gráficas, Motion for React para animación y KaTeX para fórmulas en vivo \citep{nextjs2026app,turborepo2026next,shadcn2026monorepo,reactflow2026overview,recharts2026docs,motion2026react,katex2026browser}.

\begin{table}[H]
  \centering
  \small
  \begin{tabularx}{0.96\linewidth}{l l X}
    \toprule
    \textbf{Capa} & \textbf{Tecnología} & \textbf{Responsabilidad}\\
    \midrule
    \texttt{apps/web} & Next.js + React + Recharts + KaTeX & Interfaz visual, grafo de reglas, fórmulas en vivo, descargas\\
    \texttt{packages/fuzzy-core} & TypeScript puro & Motor Mamdani determinista (única fuente de verdad)\\
    \texttt{packages/report} & LaTeX + PGFPlots + biblatex & Documento académico con figuras vectoriales y datos sincronizados\\
    \texttt{packages/presentation} & PptxGenJS & Presentación editable generada al vuelo desde el motor\\
    \bottomrule
  \end{tabularx}
  \caption{Arquitectura del monorepo. Las cuatro capas comparten una única definición del sistema difuso a través de \texttt{fuzzy-core}.}
  \label{tab:arquitectura}
\end{table}

\subsection{Motor difuso}

El paquete \texttt{fuzzy-core} contiene TypeScript puro sin dependencias en tiempo de ejecución. Sus módulos principales y su correspondencia con las etapas del modelo Mamdani (figura~\ref{fig:flujo-mamdani}) son:

\begin{itemize}
  \item \texttt{membership.ts}: implementación de las funciones triangular y trapezoidal y operación de fuzzificación.
  \item \texttt{rules.ts}: declaración de la base de reglas IF-THEN con metadatos para trazabilidad.
  \item \texttt{academic-risk-system.ts}: definición de variables, conjuntos lingüísticos y caso de prueba por defecto.
  \item \texttt{mamdani.ts}: motor de inferencia con AND mínimo, recorte, agregación máximo y construcción de la curva de salida.
  \item \texttt{centroid.ts}: defuzzificación por centroide discreto.
\end{itemize}

\subsection{Sincronización entre motor y reporte}

El script \texttt{packages/report/scripts/emit-data.ts} lee la definición del sistema directamente desde \texttt{fuzzy-core} y emite:

\begin{itemize}
  \item Archivos \texttt{.dat} con las muestras de las funciones de pertenencia (paso $0.5$).
  \item Archivos \texttt{.dat} con las curvas agregadas para cada caso de prueba.
  \item Fragmentos \texttt{.tex} con tablas de parámetros, reglas y reportes por caso.
  \item Macros LaTeX con escalas numéricas (centroide, denominador, dominante) para inserción directa en el texto.
\end{itemize}

De esta forma, cualquier modificación en el motor se refleja automáticamente en las figuras y tablas del reporte cuando se ejecuta \texttt{pnpm report}. Esto elimina el riesgo de desincronización entre código y documento.

\subsection{Interfaz web y descargas}

La interfaz expone cinco sliders, el resultado crisp, la etiqueta lingüística, el grafo de razonamiento, las funciones de pertenencia con el valor activo y las fórmulas en vivo. La barra superior contiene tres acciones de descarga:

\begin{enumerate}
  \item Presentación \texttt{.pptx} generada al vuelo con los inputs actuales del usuario.
  \item Reporte académico \texttt{.pdf} (este documento) precompilado y servido como recurso estático.
  \item Resumen ejecutivo \texttt{.pdf} generado en el navegador con \texttt{jsPDF}, conteniendo los valores actuales, gráficas vectoriales y recomendaciones.
\end{enumerate}

La captura de las gráficas vectoriales se realiza serializando el SVG renderizado por Recharts y rasterizándolo en un canvas oculto a $2\times$ de densidad, lo que garantiza nitidez en las dos exportaciones dinámicas.
